%! TEX root = main.tex

\subsection{Linear maps: what they preserve, and what they forget}

Lesson~1 described the static geometry of vector spaces. Now we study the
maps between them. A system $A\mathbf{x}=\mathbf{b}$ asks two questions:
does $\mathbf{b}$ lie among the achievable outputs of $A$, and if so, is
$\mathbf{x}$ unique? Both questions are about the map
$\mathbf{x}\mapsto A\mathbf{x}$. Differentiation is another map of the same
shape: $D(p)=p'$ is linear, and it ``forgets'' exactly the constant
polynomials. The kernel and the image make this precise.

\begin{definition}[Linear map]
A function $T\colon V\to W$ between vector spaces is \emph{linear} if
\[
  T(u+v)=T(u)+T(v), \qquad T(av)=aT(v)
\]
for all $u,v\in V$ and $a\in\R$. Equivalently,
$T(au+bv)=aT(u)+bT(v)$. By induction, $T$ preserves every linear
combination:
\[
  T(a_1v_1+\cdots+a_kv_k)=a_1T(v_1)+\cdots+a_kT(v_k).
\]
Taking $a=0$ gives $T(0)=0$: every linear map sends zero to zero.
\end{definition}

In code terms, a linear map is a pure function that distributes over the
two vector operations: the response to a combined input is the combination
of the responses (superposition). The analogy breaks where programs
branch: no thresholds, saturation, or case splits can be linear.

\begin{example}[Linear maps and an impostor]
\begin{enumerate}[label=(\alph*)]
  \item Projection $P\colon\R^3\to\R^3$, $P(x,y,z)=(x,y,0)$, is linear.
    Its outputs all lie in the $xy$-plane, and it kills exactly the
    $z$-axis. We will revisit this in the language of kernels and images.
  \item Differentiation $D$ on polynomials, $D(p)=p'$, is linear by the
    sum and constant-multiple rules for derivatives.
  \item The translation $S\colon\R^2\to\R^2$, $S(x,y)=(x+1,y)$, is
    \emph{not} linear: $S(0,0)=(1,0)\ne(0,0)$. A translation is a linear
    map plus a fixed shift (an \emph{affine} map). This is the same
    phenomenon as in Lesson~1: shifting a subspace off the origin
    destroys the structure.
\end{enumerate}
\end{example}

\begin{definition}[Kernel and image]
For a linear map $T\colon V\to W$, the \emph{kernel} and \emph{image} are
\[
  \ker T=\{v\in V:T(v)=0\}, \qquad
  \operatorname{im} T=\{T(v):v\in V\}.
\]
Note where they live: $\ker T\subseteq V$ (domain side),
$\operatorname{im}T\subseteq W$ (codomain side).
\end{definition}

\begin{theorem}[Kernels and images are subspaces]\label{thm:ker-im-subspace}
If $T\colon V\to W$ is linear, then $\ker T$ is a subspace of $V$ and
$\operatorname{im}T$ is a subspace of $W$.
\end{theorem}

\begin{proof}
Since $T(0)=0$, we have $0\in\ker T$, and $T(0)\in\operatorname{im}T$,
so both sets are nonempty. If $T(u)=T(v)=0$, then
$T(au+bv)=aT(u)+bT(v)=0$, so $\ker T$ passes the subspace test.
If $x=T(u)$ and $y=T(v)$, then $ax+by=T(au+bv)\in\operatorname{im}T$,
so $\operatorname{im}T$ passes it too.
\end{proof}

\begin{example}[Kernel and image of a matrix map]\label{ex:kernel-image-matrix}
Recall the matrix from Lesson~1:
\[
  A=\begin{pmatrix}1&2&1\\2&4&3\\1&2&2\end{pmatrix},
  \qquad T(\mathbf{x})=A\mathbf{x},\quad T\colon\R^3\to\R^3.
\]
Row reduction ($R_2-2R_1$, $R_3-R_1$, $R_3-R_2$, $R_1-R_2$) gives pivots
in columns $1$ and $3$, and $A\mathbf{x}=0$ reduces to $x_1+2x_2=0$,
$x_3=0$. Hence
\[
  \ker T=\operatorname{span}\left\{\begin{pmatrix}-2\\1\\0\end{pmatrix}\right\},
\]
which has dimension $1$. The image is the column space, spanned by the
pivot columns of the \emph{original} matrix (columns $1$ and $3$):
\[
  \operatorname{im}T
  =\operatorname{span}\left\{
    \begin{pmatrix}1\\2\\1\end{pmatrix},
    \begin{pmatrix}1\\3\\2\end{pmatrix}\right\},
\]
which has dimension $2$. Note that $1+2=3$, the dimension of the domain.
That is no coincidence.
\end{example}

\begin{theorem}[A linear map is fixed by its values on a basis]\label{thm:basis-determines-map}
Let $(v_1,\ldots,v_n)$ be a basis of $V$. For every choice of vectors
$w_1,\ldots,w_n\in W$, there is exactly one linear map $T\colon V\to W$
with $T(v_i)=w_i$ for all $i$.
\end{theorem}

\begin{proof}
Every $v\in V$ has unique coordinates $v=\sum_i a_iv_i$ (Lesson~1), so
the formula $T(v)=\sum_i a_iw_i$ is well defined. It satisfies
$T(v_i)=w_i$, and linearity follows directly from the coordinate
arithmetic. Any linear $S$ with $S(v_i)=w_i$ must satisfy
$S(v)=\sum_i a_iw_i$ by preservation of combinations, so $T$ is unique.
\end{proof}

\begin{theorem}[Maps $\R^n\to\R^m$ are matrices]\label{thm:maps-are-matrices}
Every linear map $T\colon\R^n\to\R^m$ has the form $T(\mathbf{x})=A\mathbf{x}$
for a unique $m\times n$ matrix $A$, whose $j$-th column is $T(e_j)$.
\end{theorem}

\begin{proof}
The standard basis vectors $e_1,\ldots,e_n$ form a basis of $\R^n$, so by
\Cref{thm:basis-determines-map} the values $T(e_j)$ determine $T$. Writing
$\mathbf{x}=\sum_j x_je_j$ and using linearity,
\[
  T(\mathbf{x})=\sum_j x_jT(e_j)=A\mathbf{x},
  \qquad A=\begin{pmatrix}T(e_1)&\cdots&T(e_n)\end{pmatrix}.
\]
\end{proof}

A system $A\mathbf{x}=\mathbf{b}$ therefore asks: is $\mathbf{b}$ in
$\operatorname{im}T$, and is the preimage unique?

\begin{lemma}[Exchange lemma]\label{lem:exchange}
If $(v_1,\ldots,v_n)$ spans $V$ and $(w_1,\ldots,w_m)$ is linearly
independent in $V$, then $m\le n$.
\end{lemma}

\begin{proof}
Suppose $m>n$. We show by induction on $k=0,\ldots,n$ that, after
reordering the $v$'s, the list $w_1,\ldots,w_k,v_{k+1},\ldots,v_n$ still
spans $V$. This holds for $k=0$. Assume it for some $k<n$ and write
\[
  w_{k+1}=a_1w_1+\cdots+a_kw_k+b_{k+1}v_{k+1}+\cdots+b_nv_n.
\]
Not all $b_j$ are zero: otherwise $w_{k+1}$ would be a combination of
$w_1,\ldots,w_k$, contradicting independence. Reorder
$v_{k+1},\ldots,v_n$ so that $b_{k+1}\ne0$. Then $v_{k+1}$ is a linear
combination of $w_1,\ldots,w_{k+1},v_{k+2},\ldots,v_n$ (divide by
$b_{k+1}$), so that list spans everything the old one did. At $k=n$,
$w_1,\ldots,w_n$ span $V$, so $w_{n+1}$ is their linear combination,
contradicting independence. Hence $m\le n$.
\end{proof}

\begin{corollary}[Dimension is well defined]\label{cor:dim-well-defined}
Any two finite bases of the same vector space have the same number of
vectors. This closes the open claim from Lesson~1: $\dim V$ is unambiguous.
\end{corollary}

\begin{proof}
Apply \Cref{lem:exchange} twice: each basis spans, and the other is
independent, so each size is $\le$ the other.
\end{proof}

\begin{corollary}[Extending an independent list]\label{cor:extend-basis}
In a space with a finite basis, every linearly independent list extends
to a basis.
\end{corollary}

\begin{proof}
Start with the independent list. Append the vectors of a finite basis one by one, keeping each new vector
only when it lies outside the span of everything kept so far. The kept
list stays independent by construction, and it spans $V$, since every
basis vector is either kept or was already in the span when discarded.
\end{proof}

\begin{definition}[Rank and nullity]
For a linear map $T\colon V\to W$, the \emph{rank} is $\dim\operatorname{im}T$
and the \emph{nullity} is $\dim\ker T$.
\end{definition}

\begin{theorem}[Rank--nullity]\label{thm:rank-nullity}
Let $T\colon V\to W$ be linear, with $V$ having a finite basis of $n$
vectors. Then
\[
  \dim\ker T+\dim\operatorname{im}T=n.
\]
\end{theorem}

\begin{proof}
Let $(u_1,\ldots,u_k)$ be a basis of $\ker T$, and extend it by
\Cref{cor:extend-basis} to a basis $(u_1,\ldots,u_n)$ of $V$. We claim
that $T(u_{k+1}),\ldots,T(u_n)$ is a basis of $\operatorname{im}T$.
It spans: for $v=\sum_{i=1}^n a_iu_i$,
\[
  T(v)=\sum_{i=1}^n a_iT(u_i)=\sum_{i=k+1}^n a_iT(u_i),
\]
since $T(u_i)=0$ for $i\le k$. It is independent: if
$\sum_{i=k+1}^n b_iT(u_i)=0$, then $\sum_{i=k+1}^n b_iu_i\in\ker T$, so it
equals $\sum_{j=1}^k c_ju_j$ for some coefficients $c_j$. Independence of
the full basis $(u_1,\ldots,u_n)$ forces every $b_i$ (and every $c_j$) to
be zero. Thus $\dim\operatorname{im}T=n-k$.
\end{proof}

\begin{corollary}[Injectivity test]\label{cor:injectivity-test}
$T$ is injective if and only if $\ker T=\{0\}$.
\end{corollary}

\begin{proof}
$T(u)=T(v)$ holds if and only if $T(u-v)=0$, so distinct inputs collide
exactly when a nonzero vector lies in the kernel.
\end{proof}

For a matrix map $T(\mathbf{x})=A\mathbf{x}$, rank--nullity is visible in
row reduction: with $n$ columns, each column is either a pivot column or
a free column, so
\[
  \text{(\# pivot columns)}+\text{(\# free columns)}=n,
\]
i.e.\ $\rank A+\dim\ker T=n$. The pivot \emph{indices} come from the
reduced form, but the basis of $\operatorname{im}T$ must be taken from
those columns of the \emph{original} matrix. Row operations preserve the
kernel but generally change the column space, so reading an image basis
off the reduced matrix is a trap.

\begin{remark}[Reading $A\mathbf{x}=\mathbf{b}$ through rank--nullity]
A solution exists if and only if $\mathbf{b}\in\operatorname{im}T$. If
$\mathbf{x}_p$ is one solution, every solution is
$\mathbf{x}_p+\mathbf{z}$ with $\mathbf{z}\in\ker T$ (indeed,
$A\mathbf{x}=\mathbf{b}=A\mathbf{x}_p$ gives $A(\mathbf{x}-\mathbf{x}_p)=0$): the solution set is
either empty or a shifted copy of the kernel. It is a single point
exactly when the nullity is zero.
\end{remark}

\begin{remark}[Common traps]
\begin{itemize}
  \item The kernel lives in the domain, the image in the codomain. For
    $T\colon\R^3\to\R^2$, $\ker T\subseteq\R^3$ while
    $\operatorname{im}T\subseteq\R^2$; writing them in the same space is
    a type error.
  \item ``Injective $\iff$ surjective'' needs finite dimensions and
    $\dim V=\dim W$. It fails in infinite dimensions: differentiation on
    all polynomials is surjective but kills constants, while
    multiplication by $x$ is injective but never produces a nonzero
    constant. Rank--nullity is the finite-dimensional reason the two
    coincide.
\end{itemize}
\end{remark}

In code terms, the kernel measures information loss: a nonzero kernel
means distinct inputs collide, so the input cannot be reconstructed from
the output. The rank counts the effective degrees of freedom in the
output. Real neural networks interleave linear layers with nonlinearities
precisely to escape these rigid limits.

The next step adds eigenvalues: for a map $T\colon V\to V$, the
\emph{eigenspace} for $\lambda$ is $\ker(T-\lambda I)$, and
rank--nullity will tell us how large it can be.
